\documentclass[10pt,a4paper]{amsart}
\usepackage[margin=19mm]{geometry}
\usepackage{amsmath,amssymb,mathtools}
\usepackage[scr=rsfs,cal=euler]{mathalfa}
\usepackage{xfrac}
\usepackage[unicode,hidelinks,bookmarks=false]{hyperref}
\newcommand{\ie}{i.e.\ }
\newcommand{\cf}{cf.\ }
\newcommand{\resp}{respectively\ }
\newcommand{\Subsection}{\subsection}
\providecommand{\nobreakdash}{\nobreak\hbox{-}}
\setlength{\emergencystretch}{2em}
\multlinegap0pt
\usepackage{bm}
\usepackage{bbm}
\usepackage{dsfont}
\usepackage{xspace}
\usepackage{cancel}
\usepackage{mathtools}
\usepackage{amsthm}
\makeatletter
\@ifundefined{theorem}{\newtheorem{theorem}{Theorem}}{}
\makeatother
\title[Reflection Positivity in Free Fermionic Theories]{Reflection Positivity in Free Fermionic Theories}
\author{Carl Handrack, Manfred Salmhofer}
\date{Source: arXiv:2608.10685v1 (CC BY 4.0)}
\begin{document}
\maketitle

\noindent\emph{Source and licence.} Reflection Positivity in Free Fermionic Theories. Version arXiv:2608.10685v1 (CC BY 4.0), distributed under the Creative Commons Attribution 4.0 International licence, as recorded in the arXiv licence metadata for this exact version and shown on the abstract page https://arxiv.org/abs/2608.10685v1. Full rights notice: licenses/arxiv-2608.10685-CC-BY-4.0.txt.\par\smallskip

\noindent\textit{Editorial scope.} Counted unit: the Theorem whose source label is \texttt{None} in the article (source0.tex lines 797--806, source label \texttt{None}), delivered together with the article's own complete original proof of it (source0.tex lines 808--897, one \texttt{proof} environment of 90 lines). No mathematics is written, repaired or re-derived: statement and proof are the article's own text line for line. Required context retained uncounted: none. Every definition, display and label of those lines is retained unchanged. Omitted from this package and not redistributed: 1--796, 807--807, 898--907. Nothing inside a retained excerpt is omitted or altered: every retained body is delivered exactly as the article's lines are written, so no omission receipt of any kind accompanies this package. Citations: the retained excerpts carry no citation command at all, so no bibliography entry of the article is transcribed and no reference is redistributed. Reference adaptation: none was needed and none was made; every reference inside the delivered unit resolves inside it, so no unresolved reference or citation is produced. Printed slips kept literal: no printed word, symbol, hypothesis or number of the counted statement or of its proof was altered. Layout and class: the article's own class is \texttt{amsart} with packages including marginnote, appendix, amssymb, mathtools, thmtools, thm-restate, slashed, tikz, biblatex, xcolor, hyperref, cleveref; the wrapper is the lane's pinned \texttt{amsart} editorial wrapper at 10pt on A4, which restates the theorem-like environments the retained text needs and adds the article's own macro definitions that the retained text needs (none), with extra wrapper packages (bm, bbm, dsfont, xspace, cancel, mathtools). The delivered numbering is the unit's own. Assets: no retained span contains a figure environment, a graphics-inclusion command, a PDF-page inclusion, a table or a TikZ picture; no image file is copied into this package, the package declares no asset dependency and no image file is redistributed with it. Licence: version arXiv:2608.10685v1 of this article is distributed under the Creative Commons Attribution 4.0 International licence, as recorded in the arXiv licence metadata for this exact version and shown on the abstract page https://arxiv.org/abs/2608.10685v1. A full rights notice is delivered at licenses/arxiv-2608.10685-CC-BY-4.0.txt. Characters: every retained excerpt is pure ASCII, so the delivered engine, XeLaTeX with its default Latin Modern text font, reports no missing character.
\begin{theorem}
  Let \(\Lambda>m\). Then \(C_{\Lambda,m}\) is not reflection positive: there
  exists \(f\in\mathcal S(\mathbb R_+^d)\) such that
  \begin{equation}
    I_{\Lambda,m}[f]
    =
    \left\langle\Theta f,C_{\Lambda,m}f\right\rangle_{L^2}
    <0.
  \end{equation}
\end{theorem}

\begin{proof}
  In Fourier space,
  \begin{equation}
    I_{\Lambda,m}[f]
    =
    \int_{\mathbb R^d}\mathrm d^d p\,
    \frac{e^{-p^2/\Lambda^2}}{p^2+m^2}
    \overline{\hat f(-p_0,\vec p)}
    \hat f(p_0,\vec p).
  \end{equation}
  This quadratic form extends continuously to \(L^2(\mathbb R_+^d)\). Since
  \(\mathcal S(\mathbb R_+^d)\) is dense in that space, it is enough to find an
  \(L^2\)-function of positive-time support for which the form is negative.

  Set
  \begin{equation}
    \omega_{\vec p}=\sqrt{|\vec p|^2+m^2}
  \end{equation}
  and make the ansatz that $f$ factors as
  \begin{equation}
    \hat f(p_0,\vec p)
    =
    \hat h(|\vec p|)\,
    \hat u\!\big(\tfrac{p_0}{\omega_{\vec p}}\big),
  \end{equation}
  where we want \(u\in L^2(\mathbb R_+)\). Substituting
  \(q=p_0/\omega_{\vec p}\), passing to spherical coordinates in
  \(\vec p\), and writing
  \begin{equation}
    \lambda(R)=\frac{\Lambda}{\sqrt{R^2+m^2}},
  \end{equation}
  gives
  \begin{equation}
    I_{\Lambda,m}[f]
    =
    S_{d-2}\int_0^\infty
    \frac{R^{d-2}\,\mathrm dR}{\sqrt{R^2+m^2}}
    e^{-R^2/\Lambda^2}|\hat h(R)|^2
    J_{\lambda(R)}[u],
  \end{equation}
  where
  \begin{equation}
    S_{d-2}
    =
    \frac{2\pi^{(d-1)/2}}{\Gamma((d-1)/2)}
  \end{equation}
  and
  \begin{equation}
    J_\lambda[u]
    =
    \int_{\mathbb R}\mathrm dq\,
    \frac{e^{-q^2/\lambda^2}}{q^2+1}
    \overline{\hat u(-q)}\hat u(q).
  \end{equation}

Choose $u(x)=(2-3x^2)e^{-x}\mathbf{1}_{(0,\infty)}(x)$. A direct calculation, using the recurrence
\begin{equation}
    2nI_{n+1}=2I_{n-1}+(2n-3)I_n,
    \qquad
    I_n=\int_{\mathbb{R}}\frac{e^{-q^2}}{(1+q^2)^n}\,dq,
\end{equation}
gives
\begin{equation} 
    J_1[u]
    =
    \frac{-50\sqrt{\pi}+21\pi e\,\operatorname{erfc}(1)}{80}
    <0,
\end{equation}
where the last inequality follows from
\begin{equation}
e\sqrt{\pi}\,\operatorname{erfc}(1)<1.
\end{equation}
Since $\Lambda>m$, define
\begin{equation}
R_*=\sqrt{\Lambda^2-m^2}>0,
\end{equation}
so that $\lambda(R_*)=1$. The map
\begin{equation}
R\longmapsto J_{\lambda(R)}[u]
\end{equation}
is continuous and negative at $R=R_*$. Hence, there exists an open interval $V\Subset(0,\infty)$ containing $R_*$ such that
\begin{equation}
J_{\lambda(R)}[u]<0
\qquad\text{for all }R\in V.
\end{equation}
Choose $0\neq\widehat{h}\in C_c^\infty(V)$. Since all other factors in the radial integral are positive, it follows that
\begin{equation}
I_{\Lambda,m}[f]<0.
\end{equation}
\end{proof}

\end{document}
